% VI/38 reader bridge -- additive, does not rewrite canonical solutions.
\paragraph{A homogeneous formula has a base locus.}
A proposed map to projective space
\([x]\mapsto[F_0(x):\cdots:F_m(x)]\) requires its homogeneous
coordinates to be simultaneous sections of the same line bundle;
in ordinary projective coordinates they are forms of one common
degree.  At a point where all \(F_i\) vanish, no target point is
defined.  For example, on \(\mathbb P_k^2\) the rule
\[
 [X:Y:Z]\dashrightarrow[X:Y]
\]
is undefined at \([0:0:1]\), the common zero of \(X\) and \(Y\).
It is a morphism on \(D_+(X)\cup D_+(Y)\), but not on all of
\(\mathbb P_k^2\).  This simple base-locus check must precede
any claim that a coordinate formula is a projective morphism.

\paragraph{Graded homomorphisms induce only the appropriate open map.}
Let \(\varphi:S\to T\) be a degree-preserving homomorphism
of positively graded rings, with irrelevant ideal \(S_+\).
Contracting a homogeneous prime \(\mathfrak q\in\operatorname{Proj}T\)
provides a point of \(\operatorname{Proj}S\) precisely when
\(\varphi^{-1}(\mathfrak q)\not\supseteq S_+\).
The natural domain is
\[
 U=\operatorname{Proj}T\setminus V_+(\varphi(S_+)T).
\]
For instance, take \(S=k[X,Y]\), \(T=k[u,v]\), and
\(\varphi(X)=u\), \(\varphi(Y)=0\).
The induced morphism is defined on \(D_+(u)\subset\mathbb P_k^1\),
where it is the constant map to \([1:0]\).  At \([0:1]\)
both pulled-back projective coordinates vanish; no image is defined
by this graded map.  A ring homomorphism alone therefore does not
guarantee a morphism between entire Proj schemes.

\paragraph{Closed immersion is an affine-local statement.}
A morphism \(i:Y\to X\) is a closed immersion when, locally on
\(X=\operatorname{Spec}A\), it identifies \(Y\) with
\(\operatorname{Spec}(A/I)\) for a suitable ideal \(I\subseteq A\).
The corresponding ring map \(A\twoheadrightarrow A/I\) is
surjective.  Accordingly, a surjective homogeneous graded-ring map
\(S\twoheadrightarrow S/I\) gives a closed immersion
\[
 \operatorname{Proj}(S/I)\hookrightarrow\operatorname{Proj}S.
\]
It is insufficient to establish only that a proposed map is
injective on closed points: a closed immersion must also recover
the correct structure sheaf, including nonreduced directions.

\paragraph{Worked embedding: the quadratic Veronese map.}
Consider
\[
 \nu_2:\mathbb P_k^1\longrightarrow\mathbb P_k^2,
       \qquad[s:t]\longmapsto[s^2:st:t^2].
\]
Let \(X,Y,Z\) denote the target homogeneous coordinates.
Because \(XZ-Y^2\) vanishes on the image, there is an
isomorphism of graded rings, where the Veronese subring is
regraded so that its quadratic generators have degree one,
\[
 k[X,Y,Z]/(XZ-Y^2)\ \xrightarrow{\ \sim\ }\ k[s^2,st,t^2]
       \subset k[s,t],
\]
where the displayed ring map is an isomorphism onto the second
Veronese subring.  The associated Proj agrees with
\(\operatorname{Proj}k[s,t]\).  For an explicit local check, on
\(X\ne0\) the conic is \(\operatorname{Spec}k[Y/X]\), because
\(Z/X=(Y/X)^2\).  Under \(\nu_2\), its coordinate \(Y/X\)
pulls back to \(t/s\), the affine coordinate on \(s\ne0\).
The corresponding coordinate-ring map is an isomorphism.
The \(Z\ne0\) chart is similar and the two charts cover the conic.
Thus \(\nu_2\) identifies \(\mathbb P_k^1\) with the
scheme-theoretic conic \(V_+(XZ-Y^2)\), proving that it is a
closed immersion.

\paragraph{Worked embedding: the Segre quadric.}
The Segre map
\[
 \sigma:\mathbb P_k^1\times_k\mathbb P_k^1\to\mathbb P_k^3,
 \quad([a:b],[c:d])\mapsto[ac:ad:bc:bd]
\]
has image in the quadric
\(Q=V_+(Z_{00}Z_{11}-Z_{01}Z_{10})\).
On the target chart \(Z_{00}\ne0\), let
\(r=Z_{10}/Z_{00}\) and \(q=Z_{01}/Z_{00}\).
The defining equation yields
\(Z_{11}/Z_{00}=rq\), so the coordinate ring is \(k[r,q]\).
The inverse image is the product chart \(a\ne0,\ c\ne0\),
where \(r=b/a\) and \(q=d/c\).  Its ring is also \(k[r,q]\),
and the induced map is an isomorphism.
The same calculation applies on each of the other three standard
charts; those charts cover \(Q\).  Hence \(\sigma\) is an
isomorphism onto the closed quadric, not just an injective map on
rational points.

\paragraph{Projective morphisms relative to a base.}
A morphism \(f:X\to Y\) is \emph{projective} if it factors as
a closed immersion \(X\hookrightarrow\mathbb P_Y^n\) followed
by the projection \(\mathbb P_Y^n\to Y\) for some \(n\).
This definition is relative: it asserts the existence of a
projective-space embedding \emph{over} \(Y\), not simply that
\(X\) is a projective scheme over a ground field.
Closed immersions and relative projective spaces are preserved by
base change, so a projective morphism remains projective after
any base change \(Y'\to Y\).  Projective morphisms are proper,
but the converse needs extra hypotheses and should not be
presented as a universal equivalence.
